%%%%%%%%%%  scan idx 105  |  folio 94  %%%%%%%%%%

\section*{4. Módulos formales de deformaciones}

Entre las categorías cofibradas naturales cuyas clases de isomorfía
de objetos no son, en general, prorepresentables, pero que admiten no
obstante un objeto formal versal, se encuentra la categoría cofibrada
de las deformaciones. Retomamos la situación y la notación de la
sección 1. Así, en toda esta sección, $\Lambda$ es un anillo local
noetheriano completo fijo, de cuerpo residual $k$, y
$\underline C_\Lambda$ es la categoría de las $\Lambda$-álgebras
locales artinianas de cuerpo residual $k$.

Sea $X_0$ un esquema fijo sobre $k$. Por deformación (infinitesimal)
de $X_0$ entendemos un par $(R,X)$, donde
$R\in\operatorname{Ob}(\underline C_\Lambda)$ y $X$ es un
$R$-esquema plano, junto con un diagrama conmutativo
\[
\begin{tikzcd}[column sep=huge,row sep=huge]
X_0 \arrow[r,"i_X"] \arrow[d] & X \arrow[d,"\mathrm{flat}"] \\
\operatorname{Spec}(k) \arrow[r] & \operatorname{Spec}(R)
\end{tikzcd}
\]
tal que el morfismo inducido
\[
 X_0\longrightarrow
 X\times_{\sisi{\operatorname{spec}(R)}{\operatorname{Spec}(R)}}\operatorname{Spec}(k)
\]
sea un isomorfismo. Como $R$ es local artiniano,
$i_X:X_0\longrightarrow X$ es necesariamente una inmersión cerrada y
un homeomorfismo entre los espacios topológicos subyacentes.

%%%%%%%%%%  scan idx 106  |  folio 95  %%%%%%%%%%

Si $(R,X)$ y $(R',X')$ son deformaciones de $X_0$ sobre $R$ y
$R'$, respectivamente, una flecha
\[
 (R',X')\longrightarrow(R,X)
\]
consiste en una aplicación $\varphi:R'\longrightarrow R$ en
$\underline C_\Lambda$ y un morfismo de esquemas
$\Phi:X\longrightarrow X'$ que encajan en un diagrama conmutativo
\[
\begin{tikzcd}[column sep=huge,row sep=large]
X \arrow[rr,"\Phi"] \arrow[ddd] && X' \arrow[ddd] \\
& X_0 \arrow[ul,"i_X"'] \arrow[ur,"i_{X'}"] \arrow[d] & \\
& \operatorname{Spec}(k) \arrow[dl] \arrow[dr] & \\
\operatorname{Spec}(R)
 \arrow[rr,"\operatorname{Spec}(\varphi)"']
&& \operatorname{Spec}(R').
\end{tikzcd}
\]
Esto define la categoría $\underline M_{X_0}$ de deformaciones
infinitesimales de $X_0$. Está provista del funtor
\[
 p:\underline M_{X_0}\longrightarrow\underline C_\Lambda,
 \qquad (R,X)\longmapsto R.
\]

Sea $(R,X)\in\operatorname{Ob}(\underline M_{X_0})$ y sea
$\varphi:R\longrightarrow S$ una aplicación en
$\underline C_\Lambda$. Existe un diagrama conmutativo
\[
\begin{tikzcd}[column sep=huge,row sep=huge]
X\times_{\operatorname{Spec}(R)}\operatorname{Spec}(S)
 \arrow[r,"p_1"] \arrow[d,"\mathrm{flat}"']
& X \arrow[d,"\mathrm{flat}"] \\
\operatorname{Spec}(S)
 \arrow[r,"\operatorname{Spec}(\varphi)"']
& \operatorname{Spec}(R),
\end{tikzcd}
\]
y la $\varphi$-flecha
\[
 (\varphi,p_1):(R,X)\longrightarrow
 \bigl(S,X\times_{\operatorname{Spec}(R)}\operatorname{Spec}(S)\bigr)
\]
es cocartesiana en $\underline M_{X_0}$. Por consiguiente,
$\underline M_{X_0}$ es una categoría cofibrada sobre
$\underline C_\Lambda$. Registramos los hechos siguientes.

%%%%%%%%%%  scan idx 107  |  folio 96  %%%%%%%%%%

\medskip
\noindent\textbf{4.1(a).}
Sea
\[
 (1_R,\Phi):(R,X_1)\longrightarrow(R,X_2)
\]
una flecha en $\underline M_{X_0}$. Así,
$\Phi:X_2\longrightarrow X_1$ es un $R$-morfismo cuya reducción
\[
 X_2\otimes_Rk\longrightarrow X_1\otimes_Rk
\]
es un isomorfismo. Como $R$ es artiniano y los $X_i$ son planos sobre
$R$, $\Phi$ es un isomorfismo. En consecuencia,
$\underline M_{X_0}$ es una categoría cofibrada en grupoides sobre
$\underline C_\Lambda$.

\medskip
\noindent\textbf{4.1(b).}
Supongamos que $X_0$ es afín, digamos
$X_0=\operatorname{Spec}(A_0)$. Entonces todo
$(R,X)\in\operatorname{Ob}(\underline M_{X_0})$ tiene también $X$
afín. En consecuencia, para todo
$R\in\operatorname{Ob}(\underline C_\Lambda)$, tomar secciones
globales identifica $\underline M_{X_0}(R)$ con la categoría
siguiente. Un objeto es una $R$-álgebra plana $A$ junto con una
aplicación de $R$-álgebras $j_A:A\longrightarrow A_0$ que induce un
isomorfismo
\[
 A\otimes_Rk\xrightarrow{\ \sim\ }A_0.
\]
Un morfismo $\Phi:A_1\longrightarrow A_2$ es una aplicación de
$R$-álgebras para la cual conmuta el diagrama
\[
\begin{tikzcd}[column sep=huge,row sep=huge]
A_1 \arrow[r,"\Phi"] \arrow[d,"j_{A_1}"']
& A_2 \arrow[d,"j_{A_2}"] \\
A_0 \arrow[r,equal] & A_0
\end{tikzcd}
\]

\medskip
\noindent\textbf{4.1(c).}
Consideremos la categoría cofibrada
$\widehat{\underline M}_{X_0}$ en grupoides sobre
$\widehat{\underline C}_\Lambda$. Por definición, para todo
$R\in\operatorname{Ob}(\widehat{\underline C}_\Lambda)$, la categoría
$\widehat{\underline M}_{X_0}(R)$ es canónicamente equivalente a la
categoría de esquemas formales $\mathfrak X$, ádicos sobre
$\operatorname{Spf}(R)$, tales que
\[
 \mathfrak X_n
 =\mathfrak X\otimes_R R/\mathfrak m^{n+1}
\]
es plano sobre $R/\mathfrak m^{n+1}$ para todo $n\geq0$. Si $X_0$ es
noetheriano, esto equivale a exigir que $\mathfrak X$ sea plano sobre
$\operatorname{Spf}(R)$. Así, en ese caso,
$\widehat{\underline M}_{X_0}$ es canónicamente equivalente sobre
$\widehat{\underline C}_\Lambda$ a la categoría de deformaciones de
$X_0$ sobre los objetos de $\widehat{\underline C}_\Lambda$, en el
sentido de los esquemas formales.

\medskip
\noindent\textbf{4.1(d).}
Sea $(R,X)\in\operatorname{Ob}(\underline M_{X_0})$, de modo que
existe un diagrama conmutativo
\[
\begin{tikzcd}[column sep=huge,row sep=huge]
X_0 \arrow[r,"i_X"] \arrow[d] & X \arrow[d,"\mathrm{flat}"] \\
\operatorname{Spec}(k) \arrow[r] & \operatorname{Spec}(R)
\end{tikzcd}
\]
cuyo morfismo inducido
$X_0\longrightarrow X\times_{\operatorname{Spec}(R)}
\operatorname{Spec}(k)$ es un isomorfismo.

%%%%%%%%%%  scan idx 108  |  folio 97  %%%%%%%%%%

Sea $U_0\subset X_0$ un subesquema abierto. Como
$i_X:X_0\longrightarrow X$ es una inmersión cerrada y un
homeomorfismo entre los espacios subyacentes, la restricción da un
diagrama conmutativo
\[
\begin{tikzcd}[column sep=huge,row sep=huge]
U_0 \arrow[r,"i_X|_{U_0}"] \arrow[d]
& X|_{U_0} \arrow[d,"\mathrm{flat}"] \\
\operatorname{Spec}(k) \arrow[r] & \operatorname{Spec}(R).
\end{tikzcd}
\]
El morfismo inducido
\[
 U_0\longrightarrow
 (X|_{U_0})\times_{\operatorname{Spec}(R)}\operatorname{Spec}(k)
 =(X\otimes_Rk)|_{U_0}
\]
es un isomorfismo, por lo que $(R,X|_{U_0})$ es una deformación de
$U_0$ sobre $R$. Además, toda flecha
$(\varphi,\Phi):(R,X)\longrightarrow(R',X')$ en
$\underline M_{X_0}$ se restringe a una flecha
\[
 (\varphi,\Phi|_{U_0}):(R,X|_{U_0})
 \longrightarrow(R',X'|_{U_0})
\]
en $\underline M_{U_0}$. Tenemos, por tanto, un funtor de restricción
\[
 \underline M_{X_0}\longrightarrow\underline M_{U_0}
\]
sobre $\underline C_\Lambda$.

También se puede localizar y completar en un punto. Sea
$x_0\in X_0$. Si $(R,X)\in\operatorname{Ob}(\underline M_{X_0})$,
entonces $\operatorname{Spec}(\mathcal O_{X,x_0})$ es una deformación
de $\operatorname{Spec}(\mathcal O_{X_0,x_0})$ sobre $R$, mientras
que $\operatorname{Spec}(\widehat{\mathcal O}_{X,x_0})$ es una
deformación de
$\operatorname{Spec}(\widehat{\mathcal O}_{X_0,x_0})$ sobre $R$; los
sombreros denotan la completación para la topología definida por las
potencias del ideal maximal. Existen, pues, funtores
\[
 \underline M_{X_0}
 \longrightarrow
 \underline M_{\operatorname{Spec}(\mathcal O_{X_0,x_0})}
 \longrightarrow
 \underline M_{\operatorname{Spec}(\widehat{\mathcal O}_{X_0,x_0})}.
\]

Demostraremos primero que $\underline M_{X_0}$ es una categoría
cofibrada homogénea sobre $\underline C_\Lambda$.
